\subsection{Polynomial coefficient certificates}
\label{sec:bernstein-certificates}

The polynomial comparisons below use rational coefficients on closed
intervals or rectangles. The following elementary criterion makes
their finite coefficient tables sufficient for the whole domain.

\begin{lemma}[Bernstein coefficient bounds]\label{lem:bernstein-certificate}
Let \(P\) be a polynomial on
\([a,b]\times[c,d]\), and write
\[
 Q(s,t)=P(a+(b-a)s,c+(d-c)t)
       =\sum_{p=0}^{n}\sum_{q=0}^{m}A_{pq}s^pt^q.
\]
Define
\begin{equation}\label{eq:bernstein-coefficients}
 B_{ij}=\sum_{p=0}^{i}\sum_{q=0}^{j}
 A_{pq}\frac{\binom{i}{p}}{\binom{n}{p}}
        \frac{\binom{j}{q}}{\binom{m}{q}},
 \qquad 0\le i\le n,\quad0\le j\le m.
\end{equation}
Then
\begin{equation}\label{eq:bernstein-convex-bound}
 \min_{i,j} B_{ij}\ \le P(x,y)\le\ \max_{i,j}B_{ij}
 \qquad ((x,y)\in[a,b]\times[c,d]).
\end{equation}
For a polynomial on one interval, take \(m=0\) and omit the
second variable. Degrees \(n,m\) may exceed the actual degrees.
The coefficients and both bounds are rational when the data are rational.
\end{lemma}
\begin{proof}
The binomial theorem gives
\[
 s^p=\sum_{i=p}^{n}
       \frac{\binom{i}{p}}{\binom{n}{p}}
       \binom{n}{i}s^i(1-s)^{n-i}.
\]
Apply this identity to both variables. It yields
\[
 Q(s,t)=\sum_{i=0}^{n}\sum_{j=0}^{m}B_{ij}
 \binom{n}{i}s^i(1-s)^{n-i}
 \binom{m}{j}t^j(1-t)^{m-j}.
\]
For \(s,t\in[0,1]\), the weights are nonnegative and sum to one.
This proves \eqref{eq:bernstein-convex-bound}, including the boundary.
\end{proof}
